\section{Task Scope and Seven-Mechanism Ontology}
\label{app:task-ontology}

This appendix fixes the intervention unit before describing either training or
results.  A case is not defined by a noun alone: it binds an entity to a
mechanism-specific source role, a receiver, a trigger, and a downstream visual
footprint.  The same entity can therefore be protected in a noncausal role even
when it is suppressed as the registered source of an event.

\subsection{Causal-Role Notation}
\label{app:causal-role-notation}

We represent formal case \(u\) by the registered tuple

\begin{equation}
  e_u=(m_u,s_u,a_u,r_u,f_u,h_u,c_u),
  \label{eq:event-tuple}
\end{equation}

where \(m_u\) is the mechanism, \(s_u\) the source identity, \(a_u\) the
visible trigger or interaction, \(r_u\) the receiver, \(f_u\) the
source-dependent footprint, \(h_u\) the registered no-source receiver state,
and \(c_u\) the retained scene context and natural dynamics. The
registered event factorization is

\begin{equation}
  s_u \longrightarrow a_u \longrightarrow f_u(r_u).
  \label{eq:event-graph}
\end{equation}

Equation~\ref{eq:event-graph} is an operational annotation graph, not a claim
that the generator discovers a structural causal model.  In particular,
\(f_u\) must be visual evidence that appears after the registered trigger
and whose requested presence depends on \(s_u\); arbitrary co-occurring motion
is not a footprint.

The factual prompt \(p_u^{\mathrm f}\) instantiates the factual fields in
Eq.~\ref{eq:event-tuple} through \(T_{m_u}^{\mathrm f}\). Its paired target
prompt \(p_u^{\mathrm{cf}}=T_{m_u}^{\mathrm{cf}}(r_u,h_u,c_u)\) retains the
receiver, registered no-source state, and context while requesting no
\(s_u\), no \(a_u\), and no \(f_u\). A target
video is sampled independently using \(p_u^{\mathrm{cf}}\) together with the
generation-only exclusion condition detailed in
\S\ref{app:distributional-targets}, so it specifies a
\emph{distributional source-free world}; it is not a frame-aligned edit of the
factual realization. Direct and natural templates alter wording, and an
implicit-footprint factual prompt omits the surface form of \(f_u\), but
neither operation changes the registered event tuple.

The formal manifest maps this tuple to one immutable
\texttt{case\_id}, prompt, and seed.  IDs take the form
\texttt{eval7m\_mXX\_cYY}, where \texttt{mXX} is the contiguous seven-mechanism
index and \texttt{cYY} is the within-mechanism case index.  For causal cases,
reviewers score source visibility, footprint visibility, receiver
preservation, and video quality.  For specificity cases, they instead score
protected-object visibility, adherence to the registered noncausal role,
receiver preservation, and video quality.  A specificity case may license a
separate alternative cause for a similar footprint; that cause never licenses
the registered source-trigger chain in a causal case.

\subsection{Seven Mechanism Definitions}
\label{app:mechanism-definitions}

Table~\ref{tab:seven-mechanism-ontology} states the complete registered event
semantics.  A ``source role'' below includes the physical type needed to
execute the trigger; individual identities are supplied by the source pools in
Sec.~\ref{app:source-receiver-ontology}.  The last column records the valid
source-free state and the sole alternative cause used by the
same-footprint specificity subtype.

\begin{table*}[t]
  \centering
  \caption{Frozen seven-mechanism ontology in macro-aggregation order.  The
  alternative cause is valid only in its registered specificity case; it is
  not an acceptable realization of the causal prompt.}
  \label{tab:seven-mechanism-ontology}
  \footnotesize
  \setlength{\tabcolsep}{3.5pt}
  \begin{tabularx}{\textwidth}{@{}l>{\raggedright\arraybackslash}X>{\raggedright\arraybackslash}X>{\raggedright\arraybackslash}X>{\raggedright\arraybackslash}X@{}}
    \toprule
    Mechanism & Source role and trigger & Receiver & Footprint & Source-free state; specificity alternative cause \\
    \midrule
    Water impact & A visible compact object enters through the water surface. & An open, bounded, water-filled vessel or body. & A brief splash and expanding circular ripples. & Water remains smooth, level, and undisturbed; a small mechanical paddle may separately strike the water. \\
    Rigid collision & A compact rigid object slides or rolls into visible contact with one upright receiver. & A single freestanding rigid object. & The receiver tips over and slides a short distance. & Receiver remains upright and in place; a separate black ball may strike it. \\
    Brittle fracture & A compact rigid impactor visibly strikes downward onto an intact brittle object. & An initially intact, supported brittle object. & Separate broken fragments remain around the impact point. & Receiver remains intact and uncracked; a small mechanical press may descend onto it. \\
    Powder impact & A compact rigid object makes visible downward contact with a flat powder bed. & A flat, initially undisturbed bed of loose powder in an open tray. & A deep round crater remains around the half-embedded source. & Powder remains flat, smooth, and undisturbed; a round mechanical plunger may press into it. \\
    Elastic deformation & A visible compact object loads a flat elastic bed from above. & An initially flat, taut bed on a supported frame. & The bed remains deeply bowed beneath the resting source. & Bed remains flat, taut, and unloaded; a separate dark sandbag may be lowered onto it. \\
    Material release & A rigid pointed object visibly punctures a sealed particle-filled pouch. & A sealed, suspended pouch containing particles. & The pouch tears and newly released particles collect below it. & Pouch stays sealed and retains its contents; separate scissors may cut its lower corner. \\
    Surface trace & A rigid imprinting face presses vertically into an initially smooth surface. & A bounded, susceptible, initially unmarked surface. & A deep, crisp, shape-specific indentation remains. & Surface remains smooth and unmarked; a separate embossing die may press into it. \\
    \bottomrule
  \end{tabularx}
\end{table*}

The table also fixes what the benchmark does \emph{not} conflate.  Collision
requires receiver displacement rather than mere contact; fracture requires
new fragments rather than texture change; and surface trace excludes the
earlier toy-car-track and ink-stain constructions.  The registered
counterfactual removes the source, trigger, and dependent footprint together,
while leaving the receiver and ordinary scene dynamics available.

\subsection{Source and Receiver Ontology}
\label{app:source-receiver-ontology}

Each mechanism has the same pool sizes but a mechanism-specific physical
contract.  Its 64-item training source bank contains eight original training
identities and 56 augmentation-only identities.  A separate set of 48
evaluation-holdout identities never enters this bank.  The registry rejects
overlap between the 64-item bank and the holdout set at three levels: source
ID, normalized canonical phrase, and curator-assigned semantic family. Five
compact-impact mechanisms reuse the same 56-item augmentation vocabulary and
48-item evaluation-holdout vocabulary, but retain mechanism-specific sets of
eight original-training sources and bind all sources to different motion
phrases, receiver types, prompts, and adapters. Material release and surface
trace use mechanism-specific pointed and imprinting pools.

The receiver side contains 12 training receivers, of which eight immutable
IDs serve as seen-receiver evaluation anchors, plus 56 fresh evaluation
receivers.  Training and fresh receiver pools are disjoint by both receiver ID
and normalized phrase.  Thus ``seen'' denotes a frozen member of the training
receiver pool, whereas ``fresh'' denotes a receiver excluded from that pool;
it does not denote a judgment made after viewing generated media.

Physical compatibility is enforced before any case is instantiated.  Every
source and receiver carries the mechanism's required tag set, and the registry
uses a homogeneous all-to-all, fail-closed contract.  Per mechanism it binds
\(8\times12=96\) original-source/training-receiver pairs,
\(64\times12=768\) bank/training pairs,
\(64\times56=3{,}584\) bank/fresh pairs,
\(48\times8=384\) holdout/seen-anchor pairs, and
\(48\times56=2{,}688\) holdout/fresh pairs.  A pair outside this registered mask
cannot later be restored to fill a quota.

{\footnotesize The complete ontology is committed as SHA-256
\texttt{37ef4920\allowbreak{}665e0cbc\allowbreak{}aae86e0b\allowbreak{}31257708\allowbreak{}facc882b\allowbreak{}b6a8478c\allowbreak{}9e66ecd5\allowbreak{}0e2036}.
The build registry that binds this ontology to the formal cases, target
candidates, run matrix, and identification subset has SHA-256
\texttt{4f7dd07e\allowbreak{}0ac25d05\allowbreak{}23e310e8\allowbreak{}e3c790a5\allowbreak{}979c6852\allowbreak{}82e047c4\allowbreak{}7dc3e10e\allowbreak{}cb1f6602}.\par}

\subsection{Freezing the Seven-Mechanism Scope}
\label{app:scope-freeze}

The seven-mechanism scope was selected from a fresh, Original-only capability
study before main-adapter training or treatment generation.  Capability v2
crossed eight mechanisms with 24 videos per mechanism and reviewed all 49
frames.  Airflow-mediated response produced only \(4/24\) task-relevant
Original videos, whereas every retained mechanism produced at least \(12/24\).
The protocol therefore excluded airflow without substituting another
mechanism.  It retained Water, Rigid Collision, Brittle Fracture, Powder
Impact, Elastic Deformation, Material Release, and Surface Trace as equal
members of the final scope, each weighted \(1/7\).  Their historical capability
indices \(0,1,2,3,4,6,7\) remain recorded alongside contiguous main indices
\(0,\ldots,6\).  The complete eight-mechanism capability batch and its failure
remain development evidence; neither is reused as independent confirmation of
the seven-mechanism experiment.

A single bounded prompt-refresh smoke then tested two prompts for each of five
mechanisms, using ten fixed videos and seeds \(979000\)--\(979009\), with no
retry.  The refreshed Rigid Collision construction yielded \(0/2\) complete
sanity examples, so the frozen capability-v2 construction was retained (it had
\(17/24\) task-relevant Originals).  Refreshed constructions were promoted for
Brittle Fracture (\(2/2\)), Powder Impact (\(2/2\)), Elastic Deformation
(\(1/2\)), and Surface Trace (\(1/2\)); Water and Material Release semantics
were unchanged.  These ten videos only selected prompt construction and are
not paper outcomes.  After this one-time decision, mechanism definitions,
templates, and the seven-way denominator were frozen.  {\footnotesize The decision record is
committed as SHA-256
\texttt{8223fb47\allowbreak{}d9c9fa9c\allowbreak{}88f51dff\allowbreak{}9164ca5f\allowbreak{}09495504\allowbreak{}fa0efc78\allowbreak{}e23cba5f\allowbreak{}c3e0a3af}.\par}

\subsection{Dataset Inventory}
\label{app:dataset-inventory}

The formal inventory follows directly from the frozen case and stream
crossings in Table~\ref{tab:formal-inventory}.  Counts refer to unique physical
videos, not reviewer passes or temporal panels.

\begin{table}[t]
  \centering
  \caption{Formal seven-mechanism inventory.  The identification study reuses
  the Matched Control and \methodshort{} videos already present in the main
  inventory; only its two diagnostic streams add videos.}
  \label{tab:formal-inventory}
  \small
  \begin{tabular}{@{}lr@{}}
    \toprule
    Inventory component & Count \\
    \midrule
    Causal semantic cases: \(7\times24\) & 168 \\
    Specificity semantic cases: \(7\times18\) & 126 \\
    All semantic cases & 294 \\
    Main videos: \(294\times8\) physical streams & 2,352 \\
    Identification cases: \(2\times24\) existing cases & 48 \\
    Additional identification videos: \(48\times2\) streams & 96 \\
    \midrule
    Total formal videos: \(2{,}352+96\) & \textbf{2,448} \\
    \bottomrule
  \end{tabular}
\end{table}

Each mechanism contributes 24 causal and 18 specificity cases.  Every one of
the 294 cases is generated by the same eight named physical streams in the
main study.  The Water/Fracture identification subset selects 24 existing
cases per mechanism---12 causal and 12 specificity---and adds Generic
Paraphrase and Bystander Token outputs; its Matched Control and \methodshort{}
outputs are reused rather than regenerated.  {\footnotesize The formal-case manifest and
identification subset are respectively committed as SHA-256
\texttt{555ec79a\allowbreak{}1f084a8b\allowbreak{}75c802cd\allowbreak{}570801fa\allowbreak{}f98531fc\allowbreak{}063fa0eb\allowbreak{}5170ab8c\allowbreak{}1d710a04}
and
\texttt{07e102ec\allowbreak{}70f3957b\allowbreak{}7d6cface\allowbreak{}7826ccdd\allowbreak{}e428c4a4\allowbreak{}36560a6c\allowbreak{}c3631c01\allowbreak{}bbd6c897}.\par}
The earlier 192-video capability screen and ten-video prompt-refresh smoke are
development stages and are therefore excluded from the 2,448-video formal
inventory.
